\subsection{Fast Beam Facility Experimental Methodology for the Numato and Cmod S7 FPGA Tests}
\label{subsec:fast_beam_method_numato_cmod}

The irradiation experiments described in this work were conducted at the Fast Beam Facility of the Ohio State University Research Reactor (OSURR) to evaluate FPGA-based digital systems under representative in-beam operating conditions. Testing was performed over two experimental days with the reactor maintained at a constant power level of 400~kW. On the first day, the gamma shutter was closed, and on the second day, the gamma shutter was open. This two-day arrangement provided two distinct irradiation conditions while preserving the same reactor power, general hardware configuration, and measurement approach. As a result, comparisons between days could be made with minimal changes to the underlying experimental setup.

The irradiation campaign simultaneously examined two different FPGA-based functions implemented on different hardware platforms within the same beam experiment. The Numato board was programmed with the digital plant and control design, which served as the application-oriented test case for evaluating the behavior of a real-time fixed-point control workload under irradiation. In parallel, four Cmod S7 boards were programmed with the FFSR test design and mounted to the back of the Numato board such that the FPGA chips were positioned in the beam line during exposure. This arrangement enabled the FFSR structures and the Numato-hosted digital plant to be irradiated concurrently under the same beam conditions while still representing distinct devices and workloads. In addition to the irradiated boards, one extra Cmod S7 board was operated on the bench near the control laptop and was not placed in the beam. This off-beam device served as a control reference, allowing behavior observed in the irradiated Cmod boards to be compared against a nominally unexposed implementation of the same design.

The physical configuration was selected to maximize experimental consistency while permitting multiple devices to be observed within a single irradiation campaign. By taping the four Cmod S7 boards to the back of the Numato board, the active FPGA packages could be co-located in the beam region while maintaining a compact and reproducible assembly. In this way, the Numato board and the irradiated Cmod boards experienced the same general beamline environment, allowing differences in observed behavior to be interpreted primarily in terms of device function, design implementation, and radiation response rather than large variations in placement or test procedure.

Electrical monitoring of the setup was performed using a Keysight E36312A programmable DC power supply. One monitored channel was assigned to the Numato board, while a second monitored channel was assigned to the USB dock associated with the experimental setup. Monitoring these two loads independently allowed the power behavior of the primary board under test to be separated from that of the supporting interface hardware. This was useful during irradiation because it reduced ambiguity when interpreting changes in measured current or power, particularly in cases where support electronics remained active for communication and logging but were not themselves the primary device of interest.

A Python-based logging script was used to configure the power supply, poll the measured voltage and current on the monitored channels at regular time intervals, calculate instantaneous power, and save the resulting measurements to time-stamped CSV files on the host computer. This produced a continuous record of voltage, current, and power for both channels throughout each irradiation run. The electrical data were then available for direct comparison with the functional behavior of the FPGA devices during beam exposure. In this way, anomalous digital behavior could be examined alongside board-level electrical behavior to determine whether a given event was associated with measurable changes in power demand, whether it was confined to functional outputs alone, or whether both effects occurred together.

The functional roles of the irradiated hardware differed intentionally. The Numato board implemented the digital plant and controller design developed in this work, thereby providing a representative embedded digital workload whose behavior could be monitored during irradiation. The irradiated Cmod S7 boards, by contrast, executed the FFSR design, which served as a continuously active digital test structure for observing logic-state evolution, upset manifestations, and possible interruptions in normal operation. Because the two workloads were exercised at the same time but on different devices, the experiment provided two complementary views of FPGA response in the beam: one associated with a control-oriented application design and one associated with a structured digital test architecture. The off-beam Cmod board further strengthened this methodology by providing a reference point for separating radiation-induced effects from ordinary drift, communication artifacts, or unrelated experimental disturbances.

To quantify the exposure conditions associated with each day of testing, reactor-power-based flux histories were generated and corresponding integrated fluence values were calculated as a function of time. The flux plots established the instantaneous irradiation environment during each run, while the fluence plots quantified cumulative exposure for the closed-shutter and open-shutter conditions. These data were used to interpret the observed FPGA behavior in terms of both real-time beam intensity and total delivered exposure. In addition, photographs of the experimental arrangement were retained to document board placement, mounting geometry, cabling, and the relationship between the irradiated devices and the supporting instrumentation. These images provide useful context for understanding the physical implementation of the test and improve the reproducibility of the methodology.

Overall, the experimental method was designed to assess multiple FPGA-based systems during a single irradiation campaign while preserving direct observability of both electrical and functional behavior. By operating the Numato digital plant design and four irradiated Cmod S7 FFSR boards concurrently at the OSURR Fast Beam Facility, maintaining an additional off-beam Cmod board as a control, varying the gamma shutter condition across two days at 400~kW reactor power, and continuously recording board and support-hardware power data, the experiment established a consistent and well-instrumented framework for evaluating FPGA response under neutron beam exposure.